\documentclass[12pt,a4paper]{article}

\usepackage{cmap}
\usepackage[T2A]{fontenc}
\usepackage[utf8]{inputenc}
\usepackage[russian]{babel}
\usepackage{amsmath,amssymb}
\usepackage[a4paper,margin=2cm]{geometry}

\newcommand{\R}{\mathbb{R}}
\newcommand{\eps}{\varepsilon}
\newcommand{\dd}{\,d}

\setlength{\parindent}{0pt}
\setlength{\parskip}{0.45em}
\sloppy

\begin{document}

\begin{center}
{\Large\bfseries Билет 9}
\end{center}

\noindent\fbox{%
  \begin{minipage}{\dimexpr\linewidth-2\fboxsep-2\fboxrule\relax}
  \normalfont
Несобственный интеграл. Критерий Коши сходимости интеграла. Интегралы от
знакопостоянных функций, признак сравнения. Интегралы от знакопеременных
функций, сходимость и абсолютная сходимость. Признаки Дирихле и Абеля
сходимости интегралов.
  \end{minipage}%
}

\section*{Краткий план}

Сначала несобственный интеграл определяется как предел собственных интегралов
при приближении к особой точке; случай нескольких особых точек сводится к
конечному числу интегралов с одной особенностью. Затем формулируется критерий
Коши: сходимость эквивалентна малости интегралов по хвостам. Для
знакопостоянных функций используется монотонность первообразной: интеграл
сходится тогда и только тогда, когда частичные интегралы ограничены; отсюда
следуют признаки сравнения. Для знакопеременных функций вводится абсолютная и
условная сходимость, доказывается, что абсолютная сходимость влечет обычную, а
затем приводятся признаки Дирихле и Абеля.

\section*{Определения}

\textbf{Несобственный интеграл с особенностью на правом конце.}
Пусть \(b\in\R\cup\{+\infty\}\), функция \(f\) определена на промежутке
\([a,b)\) и интегрируема в собственном смысле на каждом отрезке
\([a,\beta]\subset [a,b)\). Несобственным интегралом называется предел
\[
  \int_a^b f(x)\dd x
  =
  \lim_{\beta\to b-0}\int_a^\beta f(x)\dd x,
\]
если этот предел существует и конечен. В этом случае интеграл называется
сходящимся; если конечного предела нет, он называется расходящимся.
Для \(b=+\infty\) запись \(\beta\to b-0\) означает \(\beta\to+\infty\).

Интеграл с особенностью на левом конце определяется аналогично:
\[
  \int_a^b f(x)\dd x
  =
  \lim_{\alpha\to a+0}\int_\alpha^b f(x)\dd x.
\]
Если на конечном или бесконечном промежутке есть конечное число особых точек,
интеграл разбивают на конечную сумму интегралов с одной особенностью и требуют
сходимости каждого из них. Значение сходящегося интеграла равно сумме этих
значений.

\textbf{Знакопостоянные и знакопеременные функции.}
На данном промежутке функция называется знакопостоянной, если она не меняет
знак; для сходимости достаточно рассматривать случай \(f(x)\ge0\), потому что
при \(f(x)\le0\) переходят к \(-f(x)\). Функции, принимающие значения разных
знаков, называют знакопеременными.

\textbf{Абсолютная и условная сходимость.}
Несобственный интеграл
\[
  \int_a^b f(x)\dd x
\]
называется абсолютно сходящимся, если сходится интеграл
\[
  \int_a^b |f(x)|\dd x.
\]
Если \(\int_a^b f(x)\dd x\) сходится, но \(\int_a^b |f(x)|\dd x\) расходится,
то исходный интеграл называется условно сходящимся.

\section*{Теоремы}

Далее все утверждения записаны для особенности на правом конце промежутка
\([a,b)\), где \(b\in\R\cup\{+\infty\}\). Случай левого конца получается
заменой предела \(\beta\to b-0\) на \(\alpha\to a+0\).

\textbf{1. Критерий Коши сходимости несобственного интеграла.}
Пусть \(f\) интегрируема в собственном смысле на каждом отрезке
\([a,\beta]\subset[a,b)\). Тогда интеграл \(\int_a^b f(x)\dd x\) сходится
тогда и только тогда, когда
\[
  \forall\eps>0\ \exists\xi\in(a,b):\
  \forall b_1,b_2\in(\xi,b)\quad
  \left|\int_{b_1}^{b_2} f(x)\dd x\right|<\eps .
\]

\textbf{2. Критерий сходимости для неотрицательной функции.}
Пусть \(f(x)\ge0\) на \([a,b)\) и \(f\) интегрируема на каждом
\([a,\beta]\). Тогда
\[
  \int_a^b f(x)\dd x
  \quad\text{сходится}
  \quad\Longleftrightarrow\quad
  \sup_{\beta\in[a,b)}\int_a^\beta f(x)\dd x<+\infty.
\]

\textbf{3. Первый признак сравнения.}
Пусть \(f\) и \(g\) интегрируемы на каждом \([a,\beta]\) и
\[
  0\le f(x)\le g(x),\qquad x\in[a,b).
\]
Тогда из сходимости \(\int_a^b g(x)\dd x\) следует сходимость
\(\int_a^b f(x)\dd x\), а из расходимости \(\int_a^b f(x)\dd x\) следует
расходимость \(\int_a^b g(x)\dd x\).

\textbf{4. Второй признак сравнения.}
Пусть \(f,g\ge0\), обе функции интегрируемы на каждом \([a,\beta]\), и около
особой точки они эквивалентны в смысле сходимости интегралов:
существуют \(m,M>0\) и \(a_1\in[a,b)\), такие что
\[
  m g(x)\le f(x)\le M g(x),\qquad x\in[a_1,b).
\]
Тогда интегралы \(\int_a^b f(x)\dd x\) и \(\int_a^b g(x)\dd x\) сходятся или
расходятся одновременно. В частности, это верно, если
\(\lim_{x\to b-0}f(x)/g(x)=C\), где \(0<C<+\infty\).

\textbf{5. Абсолютная сходимость влечет сходимость.}
Если \(\int_a^b |f(x)|\dd x\) сходится, то сходится и
\(\int_a^b f(x)\dd x\).

\textbf{6. Признак Дирихле.}
Пусть \(f\) непрерывна, \(g\) непрерывно дифференцируема на \([a,b)\), и
выполнены условия:
\[
  \text{первообразная }F\text{ функции }f\text{ ограничена на }[a,b),
\]
\[
  \lim_{x\to b-0}g(x)=0,\qquad g'(x)\le0\quad\text{на }[a,b).
\]
Тогда несобственный интеграл
\[
  \int_a^b f(x)g(x)\dd x
\]
сходится.

\textbf{7. Признак Абеля.}
Пусть \(f\) непрерывна, \(g\) непрерывно дифференцируема на \([a,b)\), и
выполнены условия:
\[
  \int_a^b f(x)\dd x\ \text{сходится},\qquad
  g\ \text{ограничена на }[a,b),\qquad
  g'(x)\le0\quad\text{на }[a,b).
\]
Тогда несобственный интеграл
\[
  \int_a^b f(x)g(x)\dd x
\]
сходится.

\section*{Доказательства}

\textbf{Критерий Коши.}
Положим
\[
  \Phi(t)=\int_a^t f(x)\dd x,\qquad t\in[a,b).
\]
По определению несобственный интеграл \(\int_a^b f(x)\dd x\) сходится тогда и
только тогда, когда существует конечный предел \(\lim_{t\to b-0}\Phi(t)\).
По критерию Коши существования предела функции это равносильно условию
\[
  \forall\eps>0\ \exists\xi\in(a,b):\
  b_1,b_2\in(\xi,b)\Rightarrow
  |\Phi(b_2)-\Phi(b_1)|<\eps .
\]
Но
\[
  \Phi(b_2)-\Phi(b_1)=\int_{b_1}^{b_2}f(x)\dd x,
\]
поэтому получаем ровно записанный критерий.

\textbf{Критерий для неотрицательной функции.}
Пусть \(f(x)\ge0\) и
\[
  \Phi(\beta)=\int_a^\beta f(x)\dd x.
\]
Тогда \(\Phi\) неубывает на \([a,b)\). Следовательно, по теореме о пределе
монотонной функции существует конечный или бесконечный предел
\[
  \lim_{\beta\to b-0}\Phi(\beta)=\sup_{\beta\in[a,b)}\Phi(\beta).
\]
Поэтому несобственный интеграл сходится именно тогда, когда этот супремум
конечен.

\textbf{Первый признак сравнения.}
Из \(0\le f\le g\) следует
\[
  0\le\int_a^\beta f(x)\dd x\le\int_a^\beta g(x)\dd x
  \quad(\beta\in[a,b)).
\]
Если \(\int_a^b g(x)\dd x\) сходится, то частичные интегралы от \(g\)
ограничены сверху, значит ограничены и частичные интегралы от \(f\). По
критерию для неотрицательных функций \(\int_a^b f(x)\dd x\) сходится.
Утверждение о расходимости \(\int_a^b g(x)\dd x\) при расходимости
\(\int_a^b f(x)\dd x\) является контрапозицией первого утверждения.

\textbf{Второй признак сравнения.}
По принципу локализации сходимость интеграла не меняется, если заменить
\([a,b)\) на \([a_1,b)\). На этом хвосте выполнены оценки
\[
  m g(x)\le f(x)\le M g(x).
\]
Если \(\int_{a_1}^b g(x)\dd x\) сходится, то по первому признаку сравнения
сходится \(\int_{a_1}^b f(x)\dd x\). Если сходится
\(\int_{a_1}^b f(x)\dd x\), то из \(g(x)\le m^{-1}f(x)\) следует сходимость
\(\int_{a_1}^b g(x)\dd x\). Значит исходные интегралы имеют одинаковый
характер сходимости.

\textbf{Абсолютная сходимость.}
Если \(\int_a^b |f(x)|\dd x\) сходится, то по критерию Коши для любого
\(\eps>0\) найдется \(\xi\in(a,b)\), такое что при
\(b_1,b_2\in(\xi,b)\)
\[
  \int_{b_1}^{b_2}|f(x)|\dd x<\eps .
\]
Тогда
\[
  \left|\int_{b_1}^{b_2}f(x)\dd x\right|
  \le
  \int_{b_1}^{b_2}|f(x)|\dd x<\eps .
\]
Следовательно, для \(\int_a^b f(x)\dd x\) выполнено условие Коши, и этот
интеграл сходится.

\textbf{Признак Дирихле.}
Пусть \(F\) -- ограниченная первообразная \(f\):
\[
  |F(x)|\le C,\qquad x\in[a,b).
\]
Для \(\beta\in(a,b)\) интегрирование по частям дает
\[
  \int_a^\beta f(x)g(x)\dd x
  =
  F(\beta)g(\beta)-F(a)g(a)
  -
  \int_a^\beta F(x)g'(x)\dd x.
\]
Так как \(g'(x)\le0\) и \(g(x)\to0\), функция \(g\) неотрицательна на
\([a,b)\), а
\[
  \int_a^\beta |g'(x)|\dd x
  =
  \int_a^\beta -g'(x)\dd x
  =
  g(a)-g(\beta)\le g(a).
\]
Отсюда по признаку сравнения сходится
\[
  \int_a^b |F(x)g'(x)|\dd x
  \le
  C\int_a^b |g'(x)|\dd x.
\]
Значит \(\int_a^b F(x)g'(x)\dd x\) сходится абсолютно. Кроме того,
\(F(\beta)g(\beta)\to0\), потому что \(F\) ограничена и \(g(\beta)\to0\).
Переходя к пределу в формуле интегрирования по частям, получаем существование
конечного предела
\[
  \lim_{\beta\to b-0}\int_a^\beta f(x)g(x)\dd x,
\]
то есть сходимость \(\int_a^b f(x)g(x)\dd x\).

\textbf{Признак Абеля.}
Так как \(g\) монотонна и ограничена, существует конечный предел
\[
  g_0=\lim_{x\to b-0}g(x).
\]
Положим \(\widetilde g(x)=g(x)-g_0\). Тогда
\[
  \lim_{x\to b-0}\widetilde g(x)=0,\qquad
  \widetilde g'(x)=g'(x)\le0.
\]
Из сходимости \(\int_a^b f(x)\dd x\) следует, что функция
\[
  H(t)=\int_a^t f(x)\dd x
\]
имеет конечный предел при \(t\to b-0\), а значит ограничена. Поэтому по
признаку Дирихле сходится интеграл
\[
  \int_a^b f(x)\widetilde g(x)\dd x.
\]
Кроме того, по условию сходится \(g_0\int_a^b f(x)\dd x\). По линейности
несобственного интеграла сходится сумма
\[
  \int_a^b f(x)g(x)\dd x
  =
  \int_a^b f(x)\widetilde g(x)\dd x
  +
  g_0\int_a^b f(x)\dd x.
\]

\section*{Финальная устная версия}

Несобственный интеграл определяется как предел собственных интегралов при
приближении к особой точке:
\[
  \int_a^b f(x)\dd x=\lim_{\beta\to b-0}\int_a^\beta f(x)\dd x.
\]
Если предел конечен, интеграл сходится; если нет, расходится. При нескольких
особых точках интеграл разбивают на сумму интегралов с одной особенностью и
требуют сходимости всех частей.

Критерий Коши говорит, что \(\int_a^b f(x)\dd x\) сходится тогда и только
тогда, когда интегралы по достаточно далеким хвостам сколь угодно малы:
\[
  \forall\eps>0\ \exists\xi\in(a,b):\
  b_1,b_2\in(\xi,b)\Rightarrow
  \left|\int_{b_1}^{b_2}f(x)\dd x\right|<\eps .
\]
Доказательство: для \(\Phi(t)=\int_a^t f\) сходимость интеграла есть
существование конечного предела \(\Phi(t)\) при \(t\to b-0\), а это
эквивалентно критерию Коши для функции \(\Phi\).

Если \(f\ge0\), то \(\Phi(\beta)=\int_a^\beta f\) неубывает, поэтому интеграл
сходится тогда и только тогда, когда \(\Phi\) ограничена сверху. Отсюда сразу
следует признак сравнения: если \(0\le f\le g\), то сходимость интеграла от
\(g\) влечет сходимость интеграла от \(f\), а расходимость интеграла от \(f\)
влечет расходимость интеграла от \(g\). Если \(f\) и \(g\) сравнимы около
особой точки двусторонними оценками \(m g\le f\le M g\), то их интегралы
сходятся или расходятся одновременно.

Для знакопеременных функций вводят абсолютную сходимость:
\[
  \int_a^b f\ \text{сходится абсолютно, если сходится}\ \int_a^b |f|.
\]
Абсолютная сходимость влечет обычную, потому что из условия Коши для
\(\int |f|\) и неравенства
\[
  \left|\int_{b_1}^{b_2} f\right|\le\int_{b_1}^{b_2}|f|
\]
получается условие Коши для \(\int f\). Если \(\int f\) сходится, но
\(\int |f|\) расходится, то сходимость называется условной.

Признак Дирихле: если первообразная функции \(f\) ограничена, \(g\in C^1\),
\(g(x)\to0\) и \(g'(x)\le0\), то \(\int_a^b f(x)g(x)\dd x\) сходится. Главное
в доказательстве -- интегрирование по частям:
\[
  \int_a^\beta fg=F(\beta)g(\beta)-F(a)g(a)-\int_a^\beta Fg'.
\]
Первый граничный член имеет предел, а \(\int Fg'\) сходится абсолютно, потому
что \(|F|\le C\) и \(\int |g'|<\infty\).

Признак Абеля: если \(\int_a^b f(x)\dd x\) сходится, \(g\in C^1\), \(g\)
ограничена и монотонна, то \(\int_a^b f(x)g(x)\dd x\) сходится. Действительно,
у \(g\) есть конечный предел \(g_0\); записываем \(g=(g-g_0)+g_0\). К первой
части применяем признак Дирихле, потому что \(g-g_0\to0\), а ко второй --
сходимость \(\int f\).

\section*{Источники}

Иванов Г.Е. \emph{Лекции по математическому анализу. Часть 1}:
\texttt{IvGE\_1.pdf}, стр. 244--250, 252--262.

Основные роли источника: стр. 244--250 -- определение несобственного
интеграла, особенности и сведение к одной особенности; стр. 252--255 --
интегралы от неотрицательных функций и признаки сравнения; стр. 256--258 --
критерий Коши, абсолютная и условная сходимость; стр. 258--262 -- признаки
Дирихле и Абеля с доказательствами.

Также был выполнен поиск по \texttt{IvGE\_2.pdf}: прямого материала по теме
билета там не найдено; совпадения относятся к интегралам, зависящим от
параметра, поэтому ответ опирается на главу 8 первой части.

\end{document}
